\subsection{The two parameter ends}

Fix one splitting and write the complex parameter as $\xi$, so that
$F_\xi=(\sigma+\xi\alpha e)(P)$. Put
$c_*=-\min h$ and $\mathcal H=\{\operatorname{Im}\xi>c_*\}$.
The real family in the prompt is $\xi=ic$.
By an analytic extension across a boundary point we mean a proper
Hausdorff complex-analytic family on a full complex disc whose
restriction to the nonempty open part lying in $\mathcal H$ is
the prescribed family. At infinity we use the coordinate $u=1/\xi$.
This definition permits singular central spaces, multiplicities,
and additional central components.

\begin{statement}{P1A.5}{Proved}
There is no such analytic extension across any finite boundary
point of $\mathcal H$, or across infinity. The same conclusion
holds after any finite ramified change of the local boundary
coordinate. In particular neither end of the real interval
$c_*<c<+\infty$ has a proper Hausdorff analytic extension in
this sense.
\end{statement}
\begin{proof}
We show first that a purported extension would carry the marked
torus fibration on its whole punctured parameter disc, rather
than only on its original one-sided open set.
Take the reduced component containing that open family,
normalize it and resolve it by the analytic resolution theorems
used in P2.4. The comparison map is proper and is an isomorphism
on the given smooth open set. Proper images of the singular,
exceptional and critical loci in the parameter disc are analytic.
They omit a nonempty open set, hence are discrete. Shrinking
the disc retains only its centre. Thus its punctured restriction
is a proper holomorphic submersion with smooth compact
fourfold fibres.

On the resolved total space use the tenth inverse power of the
relative canonical line. Its proper direct image is coherent.
Remove torsion; on a smooth one-dimensional base the resulting
sheaf is locally free. Its generic rank is four, since on the
original open family $K^{-10}=f^*\mathcal O_{\mathbb P^1}(3)$
and the fibres of $f$ are connected. The analytic cohomology
and base-change complex of Grauert has constant matrix ranks
off a discrete set. Shrinking again makes every punctured
fibre comparison an isomorphism. No assertion of base change
at the central fibre is needed.

Evaluation of these four sections has no base points on the
original open set. The proper image of its base locus is
again discrete, so it has none over the remaining punctured
disc. It defines a proper map to $\mathbb P^3$ over that
disc. The image is a family of rational normal cubics there.
Here is the comparison in detail. Resolve the graph of the
meromorphic evaluation map over the full disc and take its
proper analytic image in $\mathbb P^3\times\Delta$.
This modification is an isomorphism wherever the evaluation
has no base points, hence on every punctured fibre now under
consideration. The image is projective over the disc. After
removing the torsion supported at the centre, its coherent
structure sheaf is flat away from a discrete subset. Its
Hilbert polynomial there is $3n+1$, since that is the
polynomial on the original nonempty open parameter set.
The singular and critical images are proper analytic subsets
of the parameter disc and omit that open set. Shrinking
removes them away from the centre. We therefore obtain a
smooth family of degree-three genus-zero curves on the
punctured disc; the nondegenerate embedding supplied by its
four independent sections is its complete cubic linear system.
This argument allows the basis of those sections to vary.

The recovered map to this smooth relative curve has connected
fibres. Indeed its Stein factorization has a finite factor
from a normal space to the normal relative curve. That finite
map has degree one on the original open parameter set, where
the given fibration has connected fibres. Its generic degree
is consequently one. It is finite and birational onto a normal
space, hence an isomorphism. This also proves connectedness
for the regular threefold fibres used below.

The critical image of the resulting map to this curve family
is proper analytic. Its components which dominate the
parameter disc are exactly the three continuations of the
original special values: any additional dominating component
would meet the original open parameter set. Their degrees
are one, because that is their degree on the original open
set. The multiplicities five and four distinguish the first
two, and the third is the cusp. They therefore give three
labelled sections on the punctured disc. Their cross-ratio
trivializes the genus-zero family as
$\mathbb P^1\times\Delta^*$, with the same coordinate as
the original fibration. Any vertical critical-image components
lie over a discrete set of parameters and have already been
removed. Hence the regular part is a proper holomorphic
submersion over
$(\mathbb P^1\setminus\{0,1,\infty\})\times\Delta^*$.

Every fibre of this last submersion is a complex three-torus.
Indeed any two points in its connected base are joined by
a path covered by finitely many coordinate discs. Along
each such disc the restricted family is a proper holomorphic
submersion. Catanese's Theorem~2.1, \emph{Deformation in the
large of some complex manifolds, I}, p.~264, says that a
deformation in the large of a complex torus of dimension
$n$ remains a complex torus of dimension $n$. The definition
on p.~261 is exactly a chain of these proper holomorphic
submersions; it imposes no Kähler assumption on the total
space. Starting on the original open set checks all its
hypotheses with $n=3$.

Pass to the universal cover of this regular base and fix the
integral $H^1$ marking by continuation from the given open
family. The holomorphic period planes of these tori give
a map to $\operatorname{Gr}(3,6)$. On the original open
set it is exactly $F_\xi$. Both this period map and the
explicit graph $F_\xi$ are holomorphic on the connected
cover, so the identity theorem makes them equal throughout.
The explicit graph remains a rank-three holomorphic plane
for every finite complex $\xi$, regardless of transversality;
at $\xi=\infty$ it has the Grassmannian limit
$\mathbb C e+\ker\alpha$. No choice of a continuation
marking changes the transversality of a real integral
period plane.

A full disc about a finite boundary point contains parameters
with $\operatorname{Im}\xi<c_*$. For any such parameter,
$-\operatorname{Im}\xi>\min h$ is attained by $h$ at a
regular base point: join a minimum point to a punctured
cusp neighbourhood, where $h$ is arbitrarily large, and
avoid the two isolated finite special points in this path.
The intermediate value theorem gives a regular point with
$h+\operatorname{Im}\xi=0$. At that point
$F_\xi\cap\overline{F_\xi}\ne0$, contradicting the period
decomposition of the just-proved complex torus fibre.
The same contradiction occurs in every punctured
neighbourhood of $u=0$: the function $\xi=1/u$ takes
values below the horizontal boundary line there. A finite
ramified reparameterization still maps an open neighbourhood
onto a neighbourhood of the boundary point, or onto a
punctured neighbourhood of infinity. It therefore retains
parameters on both sides and gives the same contradiction.
This proves both nonexistence assertions with all the
allowed central singularities retained.
\end{proof}

\begin{statement}{P1A.6}{Proved}
A holomorphic punctured-disc family whose fibres are all
members of this fourfold family has an interior parameter
limit. After a base change of degree at most two, it is
isomorphic to the pullback of the constructed family by
a holomorphic map of the full disc to $\mathcal H$, and
therefore has a proper smooth filling. Such a family
cannot approach either parameter end.
\end{statement}
\begin{proof}
At every punctured parameter, P1A.3 supplies local
miniversality and a local holomorphic classifying map to
$\mathcal H$. P1A.4 makes this map unique on each geometric
point. Thus these local maps agree on overlaps and define
one holomorphic function $\xi:\Delta^*\to\mathcal H$.
There is no unmentioned quotient by parameter translations:
P1A.4 proved their absence by the integral centralizer.

The Cayley coordinate
\[
 z=\frac{\xi-i(c_*+1)}{\xi-i(c_*-1)}
\]
has $|z|<1$. It extends holomorphically across the puncture
by the removable-singularity theorem. Its value there has
modulus strictly less than one, since a boundary value
would force $z$ to be constant by the maximum-modulus
principle. The inverse Cayley map therefore extends
$\xi$ with value in $\mathcal H$.

Local identifications with the pulled-back constructed
family differ by its automorphisms. These form the constant
holomorphic group $\mathbb C\times\mathbb Z/2$ over the
parameter half-plane: the translation field and the
half-period in P1A.4 extend over the entire family.
The finite component torsor on $\Delta^*$ is killed by
the double cover. The remaining additive holomorphic
torsor is zero because $H^1(\Delta^*,\mathcal O)=0$;
the punctured disc is Stein and its additive Cousin
problem has a solution. Thus after that cover the local
identifications glue to a global isomorphism. Pulling
back the already constructed proper smooth family along
the extended classifying map supplies the stated filling.
\end{proof}
